\documentclass[../DoAn.tex]{subfiles}
\pagenumbering{roman}
\begin{document}
	\begin{center}
		\Large{\textbf{ĐỀ TÀI TỐT NGHIỆP}}\\
	\end{center}
	\vspace{1cm}

	\noindent\textbf{Họ và tên sinh viên:} [HỌ VÀ TÊN SINH VIÊN]\\
	\textbf{Mã số sinh viên:} [MSSV]\\
	\textbf{Lớp/khóa:} [LỚP/KHÓA]\\
	\textbf{Giảng viên hướng dẫn:} [HỌ TÊN GIẢNG VIÊN HƯỚNG DẪN]

	\noindent\textbf{Tên đề tài:} Xây dựng ứng dụng Android trực quan hóa vệ tinh GNSS và phân tích chuyển động camera bằng Optical Flow.

	\noindent\textbf{Mục tiêu đề tài:} Đồ án hướng tới xây dựng một hệ thống thử nghiệm trên thiết bị Android cho phép người dùng quan sát trạng thái định vị GNSS, trực quan hóa vệ tinh trên bản đồ hai chiều và mô hình Trái Đất ba chiều, đồng thời phân tích chuyển động camera theo thời gian thực bằng các thuật toán Optical Flow. Hệ thống cũng cung cấp luồng xử lý video bằng mô hình RAFT trên máy chủ để giảm tải cho thiết bị di động trong các tác vụ suy luận nặng.

	\noindent\textbf{Nội dung thực hiện:} Các nội dung chính của đồ án bao gồm: (i) khảo sát bài toán định vị, trực quan hóa vệ tinh và phân tích chuyển động bằng ảnh; (ii) thiết kế kiến trúc ứng dụng Android kết hợp CameraX, OpenCV, GNSS API, OpenGL ES và dịch vụ xử lý video; (iii) triển khai các thuật toán KLT, Farneback, cơ chế đo đạc so sánh và pipeline RAFT; (iv) xây dựng các màn hình tương tác, lưu trữ phiên phân tích và trình phát kết quả; và (v) đánh giá hệ thống theo khả năng đáp ứng chức năng, tính ổn định và chi phí xử lý.

	\noindent\textbf{Sản phẩm dự kiến:} Sản phẩm của đồ án gồm mã nguồn ứng dụng Android \texttt{GNSSAndOpticalFlowApp}, mã nguồn máy chủ \texttt{OpticalFlowServer}, tài liệu kỹ thuật và báo cáo đồ án tốt nghiệp theo mẫu SOICT.

	\vspace{1.5cm}
	\begin{flushright}
		Giáo viên hướng dẫn\\
		\begin{tabular}{@{}c@{}}
			\textit{(Ký và ghi rõ họ tên)}
		\end{tabular}
	\end{flushright}
\end{document}
